\section{The Temptations in the Wilderness}

\begin{quotex}
[materialists] fear judgment. Because the future brings retribution for the past, people deny both the moral world order and the future in the sense of that moral world order. \flright{\textsc{Valentin Tomberg}}

One generation creates the destiny of the next. \flright{\textsc{Valentin Tomberg}}

\end{quotex}
Although the sequence differs between Matthew 4:1-11 and Luke 4:1-13, these are the three temptations:

\begin{itemize}
\item To turn stones into bread 
\item To jump off the pinnacle 
\item To rule the world 
\end{itemize}
The first item to note is the symmetry between the beginning of the Old Testament and the New Testament. Genesis sets the stage for the emergence of the First Adam and the three temptations in Eden:

\begin{itemize}
\item Listening to the serpent. The serpent was cunning and shifted the perspective of Adam and Eve from the vertical to the horizontal. 
\item Seeing fruit as a delight to the eyes. 
\item Eating of the fruit 
\end{itemize}
Adam and Eve succumbed to each of those temptations, and then experienced the corresponding effects.

First of all, Eve allowed the voice of the serpent to have equal influence with that of God. This was an act of disobedience and the result was doubt, i.e., having two minds in conflict with each other.

This put the tree in a new light, for now it seemed desirable. Doubt wants to be resolved by experience, i.e., to actually taste the fruit of the tree. That is greed, the opposite to poverty.

Eating the fruit, i.e., actually undergoing the experience is unchastity.

The New Testament begins with the stories of Jesus' origins ending with the Baptism in the Jordan. He is led to fast for 40 days in the wilderness. There, Satan tempts him, but Jesus resists each one of them. The following sections show Tomberg's analysis from each of the four works mentioned.

\paragraph{Anthroposophic Meditations on the New Testament}
In this work, there are two chapters devoted to the temptations in the wilderness. The NT starts with the human and leads to the divine. Curiously, Tomberg begins his discussion with Friedrich Nietzsche. Nietzsche experienced loneliness and isolation in their depths. That is, he knew the wilderness and its temptations, yet yielded to them.

His emptiness led him to seek the fullness of life in the instincts and the will to power. Then from the mountain top, Nietzsche came up with the idea of the Eternal Return. This means that the Earth has no future, but is condemned to repeat itself endlessly. Nietzsche claimed to have been inspired by superior forces. Yet then led him to devalue God, spirit, soul. The same temptations exist today, although few succumb with the same intensity.

Jesus in the wilderness was likewise isolated; the angels did not minister to him until afterwards. For, according to Tomberg, it is human freedom that must resist the temptations. The wilderness represents the Kali Yuga, the Dark Age. That is the temptation of humanity as a whole, which we see constantly repeated. Suprasensory experiences are rare because people have not yet decided:

\begin{itemize}
\item Whether to rule or to serve (power or obedience) 
\item Whether to possess the kingdoms of the past or to wander destitute into the future (wealth or poverty) 
\item Whether to desire miracles of knowledge (the authority of miracles or the chastity of knowledge) 
\end{itemize}
People are exposed to these temptations in diverse forms and Tomberg provides examples.

\subparagraph{First Temptation. Stones into Bread}
The tendency is to substitute a quantitative numerical value to everything qualitative and specific. But quantity is death and the property of quality, the living, cannot be reduced to it.

Tomberg considers money as an example of converting metal or paper into a “basket of bread”. The value of money is arbitrarily set while bread, which supports life, is subjected to the power of number. The prime example of this is in the USA where the major concern in elections is the GDP, as the spiritual and intellectual level of society declines precipitously.

\subparagraph{Second Temptation: retreat into the subconscious}
This temptation is to seek the source of life in the instinctual life of the subconscious. This is jumping from the pinnacle into the abyss, the domain of hidden instinctive urges. Thinking is difficult so the temptation is to expect miracles from the subconscious.

The true source of life, however, is the superconscious and the free life of thinking.

\subparagraph{Third temptation: materialism}
Materialism is the temptation to see the world as having no moral or spiritual guidance. This view grants them freedom from responsibility. However, when materialism is followed all the way through, the ruling intelligence behind matter will be seen to be the “prince of the world”.

The path of materialism leads through hallucination to insanity and from insanity to demonic possession. These terms are not meant in a clinical sense, and the symptoms are certainly noticeable in our time.

The fundamentals of materialism are force, chance, matter

\begin{itemize}
\item Blind force is the opposite of spiritual light, a denial of the Holy Spirit 
\item Blind chance is the opposite of the Logos, or Son 
\item Spiritless matter stands in contrast to the cosmic First Cause, or Father 
\end{itemize}
Force is unspiritual time, chance the lack of causality, and matter, the mechanization of life. The results are sleep, prostration, and death. The following meditation contains the opposite tendencies:

\begin{itemize}
\item Out of the Godhead is created humankind. 
\item In Christ death becomes life. 
\item In Spirit's cosmic thoughts, the soul awakens. 
\end{itemize}
\paragraph{Inner Development}
In the series of lectures published with the title \emph{Inner Development}, there is a brief discussion of the three temptations. The context is a critique of the three currents of contemporary intellectual life, viz., religion, art, science

\subparagraph{Religion}
Religion has succumbed to the temptation of reckoning with the Prince of this world. It has succumbed to the temptation to organize the world with the help of a power principle and take possession of it with the help of a centralized power organization. This is the temptation to rule the world, provided one bows down to Satan.

He then criticizes the Roman Church which, he claims, “strove to bring the world and its glories under its dominion.” However, in the Meditations, the view becomes more balanced.  He came to recognize that the Church is the Mystical Body of Christ, i.e., a real entity. It is the egregore of the Church, a factitious creation, that succumbs to the temptation, not the Mystical Body.

\subparagraph{Art}
Tomberg then turns his attention to art:

\begin{quotex}
Artistic creation is increasingly becoming a situation whereby the artist creates out of the deep dark underworld of his subconscious. 

\end{quotex}
That is because the artist succumbs to the temptation of jumping into the abyss.

\begin{quotex}
The artist leaps from the pinnacle of the temple of clear consciousness into the sphere of impulses, instincts, whence something is supposed to arise that is to be regarded as angelic revelation. 

\end{quotex}
The art of popular culture is now plagued with vulgar, scatological, and sexual innuendo. For some reason, this is considered to be a deep insight into the human condition.

\subparagraph{Science}
Science succumbs to the temptation to turn stones into bread:

\begin{quotex}
Modern science is based upon the conception that the dead mineral world can be the foundation of everything, and that everything living is only a consequence of movement in this mechanical, dead world. i.e., all bread arises out of stone. 

\end{quotex}
\subparagraph{Summary}
Tomberg related the temptations to the political forces rampant at that time. The specifics might not be so important now, but the temptations still arise in different contexts.

\begin{quotex}
World history is essentially nothing other than the continual karmic confrontation of humanity with the first, the second, the third, or all three temptations. 

\end{quotex}
\paragraph{Degeneration and Regeneration of Jurisprudence}
In this PhD thesis, written shortly after his conversion, Tomberg reformulates the temptations in the context of jurisprudence. Instead of religion, art, and science, the concern is now law, ethics, and religion.

\begin{quotex}
A view which presumes law, ethics, and religion to be a structure unity cannot avoid recognizing a kind of “fall” in the history of jurisprudence in the 19th and 20th centuries. A fall consists of succumbing step by step to the same three temptations to which Christ was exposed in the desert. 

\end{quotex}
\subparagraph{First Temptation}
\begin{quotex}
The turning away from the ideal of reason, and the turning towards the instinctual is, seen morally, nothing but the leap from the pinnacle of the temple into the abyss, hoping that there angels of God will lift the one falling. i.e., intelligence reigning in the darkness of the instinctual. 

\end{quotex}
Thought should be oriented toward the divine, but the temptation is to sink.

\begin{quotex}
The height of pure thought (pinnacle) oriented towards the divine (temple) has been left behind to find the reign of the nation's subconscious force (angels) in the instinctual (abyss). This is the path from faith to superstition. 

\end{quotex}
Just as an individual's personal instincts may be made, so too are the national instincts. Hence, there is a risk or a leap into the unknown. This temptation is to put the irrational above reason.

\subparagraph{Second Temptation}
Historically, there have been no intervention of angels to break the fall. Only the ground can break the fall, leading to the cult of materialism.

Hence, one generation creates the destiny of the next, although not necessarily a repetition of what came prior. Culture and morality were assumed to be determined by mechanical and material forces. Hence, the materialist generation reversed the order of the higher and the lower. In particular, the revolutionary movements of the 19th century were rooted in the primacy of the material. Quantity placed above quality.

This is the temptation to transform stones into bread. The transformation of inorganic and dead stones into organic bread is actually a reversal of the above and the below.

It treats culture (law, ethics, religion) as the product of the material (amoral and irrational).

\subparagraph{Third Temptation}
The materialist generation gave way to the positivistic generation:

Law is only what was laid down by a power according to its will and sanctioned by force. The good is only what leads to the set objective. And that objective is defined to be “truth”.

Matter is no longer fundamental, but force is. In man, force is actualized as the Will.

\begin{quotex}
Will is the reality in the life of a human—ultimately it creates and directs everything—including all of civilised and legal life. But for legal life this means a decisive change in its foundation: might replaces right. 

\end{quotex}
The temptation is to become the authority over all the kingdoms of the world, provided on bows down and worships Satan. This is predominantly the situation we find ourselves in today.

\paragraph{Meditations on the Tarot}
The three temptations have a prominent place in the Meditations. However, they are always put in contrast with the temptations in Eden, and in relation to the three vows of obedience, poverty, and chastity.

\begin{quotex}
The three vows are, in essence, memories of paradise, where man was united with God (obedience), where he possessed everything at once (poverty), and where his companion was at one and the same time his wife, his friend, his sister, and his mother (chastity). 

\end{quotex}
The work of redemption begins with the three temptations in the wilderness. However, this time the tempter was not the serpent. Rather, the tempter was the prince of the world (the new man, the superman, or other son of man who, if incarnated, would be the realisation of the promise of freedom made by the serpent.

The three temptations of the Son of Man in the wilderness were his experience of the directing impulses of evolution, namely the will-to-power, the “groping trial” and the transformation of the gross into the subtle. They signify at the same time the test of the three vows—the vows of obedience, chastity and poverty.

\subparagraph{Bread and stones}
The first temptation came after Jesus' forty day fast:

\begin{quotex}
\emph{Hunger} of the spirit, the soul and the body is the experience of \emph{emptiness} or poverty. It is therefore the vow of poverty which is put to the test when “the tempter came and said to him: If you are the Son of God. command these stones to become loaves of bread” (Matthew iv, 3). “Command these stones to become loaves of bread”—this is the very essence of the aspiration of humanity in the scientific epoch, namely to victory over poverty. Synthetic resins, synthetic rubber, synthetic fibre, synthetic vitamins, synthetic proteins and. . .eventually synthetic bread! —When? Soon, perhaps. Who knows? 

\end{quotex}
Forty-five years later, we can begin to answer such questions. Synthetic meat is being grown in the laboratory and will soon be mass produced. Of course, the advances in artificial intelligence are seen as a threat. Synthetic humans are already being produced for specialized purposes (i.e., sex).

Through the media, the population is being prepared to accept these androids as fully human. Some even predict them to be the next root race that will replace biological humans.

In cultural terms, this temptation leads the intellectuals to deny life, which is regarded as merely complex molecules. In the political realm, economic life is regarded as primary.

\subparagraph{Groping Trial}
This is the temptation to jump off the pinnacle into the abyss, assuming God will come to the rescue. In practical terms, it is the idea of evolution over creation. Creation is the accomplishment of absolute wisdom and absolute goodness.

Evolution, on the other hand, proceeds blindly by trial and error, from one species to the next. Evolution is actually guided by the serpent, or prince of the word. This temptation is opposed to chastity, and fornication is threefold: spiritual, psychic, and carnal.

\begin{quotex}
The principle of spiritual fornication is therefore the preference of the subconscious to the conscious and superconscious. of instinct to the Law, and of the world of the serpent to the world of the WORD. 

\end{quotex}
\subparagraph{Transformation Temptation}
The third temptation is directed against obedience. The temptation is the will-to-power, the desire to rule over the world.

\begin{quotex}
It is a matter, therefore, of accepting the ideal of the superman (“fall down and worship me”), who is the summit of evolution (“he took him to a very high mountain”) and who, having passed through the mineral, plant, animal and human kingdoms, subjecting them to his power, is lord over them, i.e. he is their final cause or aim and ideal, their representative or their collective concentrated will, and he is their master, who has taken their subsequent evolution into his hands. Now, the choice here is between the ideal of the superman, who is “as God”, and God himself. 

\end{quotex}
Obedience is faithfulness to the living God.

\subparagraph{Fundamental Law of Magic}
A fundamental law of sacred magic is this:

\begin{quotex}
That which is above being as that which is below, renunciation below sets in motion forces of accomplishment above and the renunciation of that which is above sets in motion forces of accomplishment below. 

\end{quotex}
This means:

\begin{quotex}
When you resist a temptation or renounce something desired below, you set in motion by this very fact forces of realisation of that which corresponds above to that which you come to renounce below. 

\end{quotex}
Or

\begin{quotex}
It is not desire which bears magical realisation, but rather the renunciation of desire (that you have formerly experienced, of course). For renunciation through indifference has no moral — and therefore no magical — value. 

\end{quotex}
This renunciation is actually the practice of the three sacred vows, which is the true magical training, and concentration without effort. That is the esoteric value to resisting those temptations.

\flrightit{Posted on 2018-09-07 by Cologero }
